
In order to understand the bin-to-bin migrations  due to detector resolution and similar effects a migration or smearing matrix $S$ is constructed. 
This enables a conversion of the number of true events in a true bin $j$, to the number of observed events in a reconstructed bin $i$:

\begin{equation}
\label{eq:smearing_matrix_2d}
\nu_i = \sum_{j=1}^{M} S_{ij} \mu_j
\end{equation}

where $S$ is given by:

\begin{equation}
S_{ij} = P(\text{observed in bin}\, i \,|\, \text{true value in bin}\, j)
\end{equation}

and $\nu_i$ and $\mu_j$ are the numbers of events in the reconstructed bin $i$ and true bin $j$ respectively, and $M$ is the total number of bins.  The matrix is determined by Monte Carlo techniques using the model discussed in Sec.~\ref{sec:evgen} and detector simulation discussed in Sec.~\ref{sec:detsim}.

Since data are measured in reconstructed variables only, data can be considered 'smeared'. There are a number of ways this data can be analyzed to allow direct comparison with theoretical calculations.
Data can be unsmeared, i.e. corrected for resolution through unfolding~\cite{DAgostini:1994fjx}. 
This is used in experiments such as \minerva~\cite{Ruterbories:2018gub}.  
However, Unfolding have uncertainties when using the finite samples from Monte Carlo samples.
T2K used unfolding in early papers and has started to consider forward folding techniques~\cite{Koch2019}.  

 Using unfolding techniques, the migration matrix is used to deconvolve the resolution effects from the data.  This allows direct comparison between data and theory and between experiments.
 In forward folding, a migration matrix is built with Monte Carlo techniques.  Although the data are presented without modification, any theory distribution has to be smeared for a direct comparison.  In addition, data from different experiments can't be directly compared.  

This analysis uses forward folding, which avoids the inherent problems that arise with unfolding techniques.
To enable comparison between theoretical models and data the smearing matrix will be published as Supplementary Data. Using this forward-folding technique we define:

\begin{equation}
\label{eq:xsec_smeared}
\begin{split}
\left ( \frac{d\sigma}{dp_\mu} \right )_i &= \frac{N_i - B_i}{ \tilde{\epsilon}_i \cdot N_{target} \cdot \Phi_{\nu_\mu} \cdot (\Delta p_\mu)_i}\\
%\left ( \frac{d\sigma}{d\cos\theta_\mu} \right )_i &= \frac{N_i - B_i}{\tilde{\epsilon}_i \cdot N_{target} \cdot \Phi_{\nu_\mu} \cdot (\Delta \cos\theta_\mu)_i}\\
%\frac{d^2\sigma}{dp_{\mu, i} \, d\cos\theta_{\mu, j}} &= \frac{N_{ij} - B_{ij}}{\tilde{\epsilon}_{ij} \cdot N_{target} \cdot \Phi_{\nu_\mu} \cdot (\Delta p_\mu)_i \cdot (\Delta \cos\theta_\mu)_j} \\
\end{split}
\end{equation}

where $\tilde{\epsilon}_i$ is the new efficiency, now as a function of the reconstructed quantities:

\begin{equation}
\label{eq:eff_smear}
\tilde{\epsilon}_i = \frac{ \sum_{j=1}^{M} S_{ij}N^\text{sel}_j}{ \sum_{j=1}^{M} S_{ij}N^\text{gen}_j} \quad 
\end{equation}

where $N^\text{sel}_j$ is the number of signal selected events in true bin $j$, and $N^\text{gen}_j$ is the number of generated signal events in true bin $j$.

